\documentclass[12pt,a4paper]{article}

% --- Metadata (per course) ---
\def\courseshortheader{XSTK NÂNG CAO}
\def\coverbookline{NHÓM 5 --- XÁC SUẤT VÀ THỐNG KÊ}
\def\covermaintitle{Xác suất thống kê nâng cao}
\def\coversubtitle{Giáo án lý thuyết --- Ước lượng, kiểm định, thống kê tiệm cận}

% Shared preamble for nhom_5 (requires \courseshortheader before \input)
\usepackage[utf8]{inputenc}
\usepackage[vietnamese]{babel}
\usepackage{amsmath,amssymb,amsthm}
\usepackage{geometry}
\usepackage{enumitem}
\usepackage[most]{tcolorbox}
\usepackage{hyperref}
\usepackage{fancyhdr}
\usepackage{graphicx}
\usepackage{booktabs}
\usepackage{tikz}
\usetikzlibrary{calc,arrows.meta,decorations.pathreplacing,positioning}
\usepackage{eso-pic}
\usepackage{xcolor}

\geometry{a4paper, margin=2cm, bottom=2.5cm}
\hypersetup{colorlinks=true, linkcolor=blue!70!black, urlcolor=blue}
\setlist[itemize]{topsep=4pt, itemsep=2pt}
\setlist[enumerate]{topsep=4pt, itemsep=2pt}

\AddToShipoutPictureFG{%
  \AtPageCenter{%
    \begin{tikzpicture}[remember picture, overlay]
      \node[rotate=45, scale=1, opacity=0.06]{%
        \fontsize{60}{72}\selectfont\bfseries\sffamily Đặng Tấn Quí%
      };
    \end{tikzpicture}%
  }%
}

\pagestyle{fancy}
\fancyhf{}
\fancyhead[L]{\small\textit{\courseshortheader}}
\fancyhead[R]{\small\textit{Đặng Tấn Quí}}
\fancyfoot[C]{\thepage}
\fancyfoot[L]{\footnotesize\textcopyright\ Đặng Tấn Quí}
\fancyfoot[R]{\footnotesize SĐT: 0377\,122\,210}
\renewcommand{\headrulewidth}{0.4pt}
\renewcommand{\footrulewidth}{0.4pt}

\fancypagestyle{plain}{%
  \fancyhf{}
  \fancyfoot[C]{\thepage}
  \fancyfoot[L]{\footnotesize\textcopyright\ Đặng Tấn Quí}
  \fancyfoot[R]{\footnotesize SĐT: 0377\,122\,210}
  \renewcommand{\headrulewidth}{0pt}
  \renewcommand{\footrulewidth}{0.4pt}
}

\newtcolorbox{dongco}[1][]{%
  enhanced, breakable,
  colback=teal!4!white, colframe=teal!60!black,
  fonttitle=\bfseries, title={Động cơ học -- Tại sao cần khái niệm này?}, #1%
}
\newtcolorbox{ynghia}[1][]{%
  enhanced, breakable,
  colback=purple!3!white, colframe=purple!50!black,
  fonttitle=\bfseries, title={Ý nghĩa \& Trực giác}, #1%
}
\newtcolorbox{ungdung}[1][]{%
  enhanced, breakable,
  colback=olive!5!white, colframe=olive!60!black,
  fonttitle=\bfseries, title={Ứng dụng thực tế}, #1%
}
\newtcolorbox{dinhnghia}[1][]{%
  enhanced, breakable,
  colback=blue!3!white, colframe=blue!60!black,
  fonttitle=\bfseries, title={Định nghĩa}, #1%
}
\newtcolorbox{dinhly}[1][]{%
  enhanced, breakable,
  colback=green!3!white, colframe=green!50!black,
  fonttitle=\bfseries, title={Định lý}, #1%
}
\newtcolorbox{tinhchat}[1][]{%
  enhanced, breakable,
  colback=yellow!5!white, colframe=orange!50!black,
  fonttitle=\bfseries, title={Tính chất}, #1%
}
\newtcolorbox{phuongphap}[1][]{%
  enhanced, breakable,
  colback=gray!5!white, colframe=gray!50!black,
  fonttitle=\bfseries, title={Phương pháp}, #1%
}
\newtcolorbox{luuy}[1][]{%
  enhanced, breakable,
  colback=red!3!white, colframe=red!50!black,
  fonttitle=\bfseries, title={Lưu ý}, #1%
}
\newtcolorbox{congthuc}[1][]{%
  enhanced, breakable,
  colback=cyan!3!white, colframe=cyan!50!black,
  fonttitle=\bfseries, title={Công thức}, #1%
}
\newtcolorbox{vidu}[1][]{%
  enhanced, breakable,
  colback=violet!3!white, colframe=violet!50!black,
  fonttitle=\bfseries, title={Ví dụ}, #1%
}
\newtcolorbox{phanvidu}[1][]{%
  enhanced, breakable,
  colback=red!2!white, colframe=red!40!black,
  fonttitle=\bfseries, title={Phản ví dụ}, #1%
}
\newtcolorbox{tongket}[1][]{%
  enhanced, breakable,
  colback=blue!4!white, colframe=blue!70!black,
  fonttitle=\bfseries, title={Tổng kết chương}, #1%
}


\begin{document}

\thispagestyle{empty}
\begin{tikzpicture}[remember picture, overlay]
  \fill[blue!80!black] (current page.south west) rectangle (current page.north east);
  \fill[blue!60!cyan, opacity=0.8]
    (current page.south west) -- (current page.north west) --
    ([xshift=-3cm]current page.north east) -- (current page.south east) -- cycle;
  \foreach \r/\o in {6/0.05, 4.5/0.08, 3/0.12} {
    \fill[white, opacity=\o] ([xshift=2cm, yshift=-2cm]current page.north east) circle (\r cm);
  }
  \draw[white, line width=2pt] ([yshift=-5cm]current page.north west) ++(2cm,0) -- ++(17cm,0);
  \draw[white, line width=2pt] ([yshift=-16cm]current page.north west) ++(2cm,0) -- ++(17cm,0);
  \node[anchor=west, white, font=\fontsize{28}{34}\bfseries\selectfont]
    at ([xshift=3cm, yshift=-7.5cm]current page.north west) {\coverbookline};
  \node[anchor=west, white, font=\fontsize{20}{26}\selectfont]
    at ([xshift=3cm, yshift=-9cm]current page.north west) {\covermaintitle};
  \node[anchor=west, white, font=\fontsize{16}{22}\selectfont]
    at ([xshift=3cm, yshift=-10.2cm]current page.north west) {\coversubtitle};
  \node[anchor=west, white, font=\Large\bfseries]
    at ([xshift=3cm, yshift=-13cm]current page.north west) {Biên soạn:};
  \node[anchor=west, yellow!80!white, font=\fontsize{24}{30}\bfseries\selectfont]
    at ([xshift=3cm, yshift=-14.5cm]current page.north west) {ĐẶNG TẤN QUÍ};
  \node[anchor=west, white!90, font=\large]
    at ([xshift=3cm, yshift=-18cm]current page.north west) {SĐT: 0377\,122\,210};
  \node[anchor=south, white!80, font=\small]
    at ([yshift=1.5cm]current page.south) {Tài liệu có bản quyền --- Cấm sao chép khi chưa được sự đồng ý của tác giả};
  \node[anchor=south east, white, font=\Large\bfseries]
    at ([xshift=-2cm, yshift=3cm]current page.south east) {\the\year};
\end{tikzpicture}

\newpage
\tableofcontents
\newpage

\section*{Lời nói đầu}
\addcontentsline{toc}{section}{Lời nói đầu}

Tài liệu thuộc \textbf{Nhóm 5 --- Xác suất và thống kê}, biên soạn theo cùng triết lý với giáo án Đại số tuyến tính: học để \emph{hiểu} và \emph{dùng}, không chỉ để ghi nhớ công thức.

\begin{enumerate}
  \item \textbf{Động cơ trước định nghĩa.} Mỗi phần lớn mở bằng ô \textsf{\color{teal!60!black} Động cơ học} --- câu hỏi thực tế hoặc bài toán mẫu cần khái niệm đó.
  \item \textbf{Trực giác trước công thức.} Ô \textsf{\color{purple!50!black} Ý nghĩa \& Trực giác} giải thích ``công thức đang nói gì'' (xác suất = tần suất dài hạn, kỳ vọng = trung bình có trọng số, v.v.).
  \item \textbf{Ví dụ có kiểm tra.} Ô \textsf{\color{violet!50!black} Ví dụ} nên tự làm trước khi đọc lời giải; phần thống kê ứng dụng cần gắn dữ liệu số cụ thể.
  \item \textbf{Tổng kết cuối mỗi chương.} Ô \textsf{\color{blue!70!black} Tổng kết chương} liệt kê kết quả phải nhớ và liên kết sang phần sau.
  \item \textbf{Phân tầng theo môn.} X1--X3 là nền; X4--X5 nâng cao lý thuyết; X6--X8 chuyên sâu (đa biến, Bayes, quá trình ngẫu nhiên).
\end{enumerate}

\noindent\textbf{Cách đọc:} Động cơ $\to$ Định nghĩa / Định lý $\to$ Ví dụ $\to$ Ý nghĩa $\to$ Tổng kết. Với môn thống kê, luôn hỏi: \emph{mẫu là gì, tham số nào chưa biết, giả thuyết $H_0$ là gì}.

\newpage


\section{Mở đầu}

\begin{dongco}[title={Động cơ --- Mở đầu}]
Mở rộng X2: tính chất ước lượng (hiệu quả Cramér--Rao), kiểm định likelihood ratio, delta method và phân phối tiệm cận --- nền cho nghiên cứu và học sau đại học.
\end{dongco}

\newpage

\section{Lý thuyết ước lượng nâng cao}

\begin{dongco}[title={Động cơ --- Lý thuyết ước lượng nâng cao}]
Khi nào ước lượng ``tốt nhất'' và sai số có cận trên không?
\end{dongco}

\begin{dinhnghia}[title={Họ mũ (exponential family)}]
  $f(x;\theta) = h(x)\,c(\theta)\,\exp\!\left(\sum_{j=1}^k \eta_j(\theta)\,T_j(x)\right)$.

  Thống kê đủ tối thiểu: $\bigl(T_1(\mathbf{X}), \ldots, T_k(\mathbf{X})\bigr)$.
\end{dinhnghia}

\begin{dinhly}[title={Lehmann--Scheffé}]
  Nếu $T$ đủ đầy đủ (complete sufficient) và $g(T)$ không chệch cho $\tau(\theta)$ thì $g(T)$ là ước lượng không chệch phương sai nhỏ nhất duy nhất (UMVUE).
\end{dinhly}

\begin{dinhnghia}[title={Tính hiệu quả tiệm cận}]
  $\hat{\theta}_n$ \textbf{tiệm cận hiệu quả} nếu $\sqrt{n}(\hat{\theta}_n - \theta) \xrightarrow{d} \mathcal{N}(0, I(\theta)^{-1})$.

  MLE là tiệm cận hiệu quả dưới điều kiện chính quy.
\end{dinhnghia}

\begin{phuongphap}[title={Phương pháp delta}]
  Nếu $\sqrt{n}(\hat{\theta} - \theta) \xrightarrow{d} \mathcal{N}(0, \sigma^2)$ và $g$ khả vi, $g'(\theta) \ne 0$:
  \[
    \sqrt{n}\bigl(g(\hat{\theta}) - g(\theta)\bigr) \xrightarrow{d} \mathcal{N}\!\left(0, [g'(\theta)]^2\sigma^2\right).
  \]
\end{phuongphap}

%% ================================================================
%%  CHƯƠNG 2: KIỂM ĐỊNH NÂNG CAO
%% ================================================================

\begin{tongket}[title={Tổng kết: Lý thuyết ước lượng nâng cao}]
\begin{enumerate}
  \item Nắm lại các \textbf{định nghĩa} và \textbf{công thức} cốt lõi trong mục ``Lý thuyết ước lượng nâng cao''.
  \item Tự làm lại ít nhất một ví dụ (nếu có) hoặc áp dụng công thức vào một tình huống số.
  \item Ghi chú điều kiện áp dụng (giả thiết độc lập, phân phối chuẩn, $n$ đủ lớn, v.v.).
\end{enumerate}
\end{tongket}

\section{Kiểm định nâng cao}

\begin{dongco}[title={Động cơ --- Kiểm định nâng cao}]
So sánh các quy tắc kiểm định: Neyman--Pearson, LRT, Wald.
\end{dongco}

\begin{dinhly}[title={Kiểm định tỉ số hợp lý tổng quát (GLRT)}]
  $\Lambda = \dfrac{\sup_{\theta \in \Theta_0} L(\theta)}{\sup_{\theta \in \Theta} L(\theta)}$. Bác bỏ $H_0$ khi $\Lambda < c$.

  Dưới $H_0$ và điều kiện chính quy: $-2\ln\Lambda \xrightarrow{d} \chi^2(r)$, $r = \dim\Theta - \dim\Theta_0$.
\end{dinhly}

\begin{dinhly}[title={Kiểm định Wald}]
  $W = \dfrac{(\hat{\theta} - \theta_0)^2}{\widehat{\text{Var}}(\hat{\theta})} \xrightarrow{d} \chi^2(1)$ dưới $H_0: \theta = \theta_0$.
\end{dinhly}

\begin{dinhly}[title={Kiểm định Score (Rao)}]
  $S = \dfrac{U(\theta_0)^2}{I(\theta_0)} \xrightarrow{d} \chi^2(1)$, \;$U(\theta) = \dfrac{\partial \ell}{\partial \theta}$.

  Ưu điểm: chỉ cần tính tại $\theta_0$ (không cần MLE).
\end{dinhly}

\begin{luuy}
  Ba kiểm định GLRT, Wald, Score tiệm cận tương đương. Mẫu hữu hạn: GLRT thường tốt hơn.
\end{luuy}

%% ================================================================
%%  CHƯƠNG 3: THỐNG KÊ TIỆM CẬN
%% ================================================================

\begin{tongket}[title={Tổng kết: Kiểm định nâng cao}]
\begin{enumerate}
  \item Nắm lại các \textbf{định nghĩa} và \textbf{công thức} cốt lõi trong mục ``Kiểm định nâng cao''.
  \item Tự làm lại ít nhất một ví dụ (nếu có) hoặc áp dụng công thức vào một tình huống số.
  \item Ghi chú điều kiện áp dụng (giả thiết độc lập, phân phối chuẩn, $n$ đủ lớn, v.v.).
\end{enumerate}
\end{tongket}

\section{Thống kê tiệm cận}

\begin{dongco}[title={Động cơ --- Thống kê tiệm cận}]
Khi $n$ lớn, phân phối của ước lượng xấp xỉ chuẩn --- cơ sở khoảng tin và kiểm định xấp xỉ.
\end{dongco}

\begin{dinhnghia}[title={Hiệu quả tương đối tiệm cận (ARE)}]
  $\text{ARE}(\hat{\theta}_1, \hat{\theta}_2) = \dfrac{\text{Var}_{\text{asy}}(\hat{\theta}_2)}{\text{Var}_{\text{asy}}(\hat{\theta}_1)}$.

  Số mẫu cần thiết để $\hat{\theta}_2$ đạt cùng độ chính xác với $\hat{\theta}_1$.
\end{dinhnghia}

\begin{dinhly}[title={Định lý Glivenko--Cantelli}]
  $\sup_x |\hat{F}_n(x) - F(x)| \xrightarrow{\text{a.s.}} 0$ (hàm phân phối thực nghiệm hội tụ đều a.s.).
\end{dinhly}

\begin{dinhly}[title={Hệ quả CLT cho $\hat{F}_n$}]
  $\sqrt{n}\bigl(\hat{F}_n(x) - F(x)\bigr) \xrightarrow{d} \mathcal{N}\!\left(0, F(x)(1-F(x))\right)$.
\end{dinhly}

\begin{phuongphap}[title={Bootstrap}]
  \begin{enumerate}
    \item Rút mẫu lặp $B$ lần từ $\hat{F}_n$ (hoặc tham số từ $\hat{\theta}$).
    \item Tính thống kê $T^*_b$ trên mỗi mẫu bootstrap.
    \item Xấp xỉ phân phối, CI, sai số chuẩn từ $(T^*_1, \ldots, T^*_B)$.
  \end{enumerate}
\end{phuongphap}

\begin{phuongphap}[title={Phương pháp EM (Expectation--Maximization)}]
  Cho dữ liệu không đầy đủ $(X, Z)$, $Z$ ẩn:
  \begin{enumerate}
    \item \textbf{E-step:} $Q(\theta|\theta^{(t)}) = E_{Z|X,\theta^{(t)}}[\ln L(\theta; X, Z)]$.
    \item \textbf{M-step:} $\theta^{(t+1)} = \arg\max_\theta Q(\theta|\theta^{(t)})$.
  \end{enumerate}
  $\ell(\theta^{(t+1)}) \ge \ell(\theta^{(t)})$ (tăng đơn điệu). Hội tụ về điểm dừng (cực đại địa phương).
\end{phuongphap}

\begin{tongket}[title={Tổng kết: Thống kê tiệm cận}]
\begin{enumerate}
  \item Nắm lại các \textbf{định nghĩa} và \textbf{công thức} cốt lõi trong mục ``Thống kê tiệm cận''.
  \item Tự làm lại ít nhất một ví dụ (nếu có) hoặc áp dụng công thức vào một tình huống số.
  \item Ghi chú điều kiện áp dụng (giả thiết độc lập, phân phối chuẩn, $n$ đủ lớn, v.v.).
\end{enumerate}
\end{tongket}



\end{document}
